%% EvoRank short paper, CEURART one-column format.
%% Compile with pdflatex or on Overleaf with ceurart.cls alongside.
\documentclass[]{ceurart}
\sloppy
\usepackage{booktabs}

\begin{document}

\copyrightyear{2026}
\copyrightclause{Copyright for this paper by its authors.
  Use permitted under Creative Commons License Attribution 4.0 International (CC BY 4.0).}

%% TODO: replace with the exact workshop line provided by the organizers.
\conference{Workshop information to be inserted upon acceptance}

\title{When Does LLM-Guided Discovery Actually Discover? An Audited Two-Campaign Study on E-Commerce Learning-to-Rank}

\author[1]{Shabaz Patel}[%
email=rayhan.b.patel@gmail.com,
]
\cormark[1]
\author[1]{Rayhan Patel}[%
]
\address[1]{Independent Researcher}
\cortext[1]{Corresponding author.}
%% TODO: add Rayhan's email and any ORCID iDs; adjust corresponding author if needed.

\begin{abstract}
LLM-guided evolutionary loops are increasingly used to discover machine
learning components, and their results are typically reported on the fold the
loop itself selected on. We present an open, fully audited case study on
multi-objective e-commerce Learning-to-Rank that asks when such loops
discover anything real. In a first campaign we evolve LambdaMART-style
training objectives over relevance, conversion, and revenue signals on the
Expedia ICDM 2013 dataset. The loop appears to work, beating baselines on its
selection fold, yet a transfer audit on 60k held-out queries shows the gains
are almost entirely fitness noise: per-edit effect sizes (0.002 to 0.005
NDCG@10) are smaller than the selection-fold noise floor, and neither seeded
domain knowledge nor gradient-introspective diagnostic feedback changes what
transfers, although seeded knowledge accelerates selection-fold convergence
about five-fold. In a second campaign we point the same loop at a search
space with measured headroom above the noise floor: feature construction,
model selection, loss selection, and ensemble architecture. There, three
independent seeds converge within 50 iterations (about ten dollars each) on
interpretable pipelines that Pareto-dominate an Optuna-tuned LambdaMART on
all three objectives on held-out data, an advantage that persists under
config-equalized comparison when training on the full 6.9M-impression
dataset. Our central finding is diagnostic rather than architectural:
search-space headroom relative to fitness noise, not guidance sophistication,
separates apparent discovery from real discovery. We release the loop, the
audit stack, and a measured catalog of failure modes, including three
generations of target-encoding leakage, each caught by the audit.
\end{abstract}

\begin{keywords}
  learning to rank \sep
  LLM agents \sep
  automated machine learning \sep
  evolutionary program search \sep
  e-commerce search \sep
  reproducibility
\end{keywords}

\maketitle

\section{Introduction}
Automated discovery loops that pair a large language model with evolutionary
search over programs have produced striking results, and industrial systems
such as Meta's Ranking Engineer Agent apply the pattern to ranking at scale.
These systems are closed, and published results in the genre share a
methodological gap: gains are overwhelmingly reported on the data the loop
used to select candidates. For ranking, where realistic per-candidate
evaluations are noisy by necessity, this gap matters. Selection under noise
systematically promotes lucky candidates, and a loop can appear to make
steady progress while discovering nothing that transfers.

We contribute an open, end-to-end audited account of when the pattern works
and when it silently does not, on a public e-commerce ranking task with three
competing objectives (relevance, conversion, revenue). Our contributions are:
(1) a measured failure: evolved training objectives whose selection-fold
gains do not survive a held-out transfer audit, with the mechanism quantified
(effect sizes below the fitness noise floor); (2) guidance ablations showing
that seeded domain knowledge and evidence-rich diagnostic feedback change the
speed and style of search but not what generalizes; (3) a measured success
under identical auditing: pointing the same loop at a search space whose
effect sizes clear the noise floor yields transfer-audited, Pareto-dominant
pipelines in all three seeds, which independently reconstruct the
load-bearing structure of the ICDM 2013 challenge winners' hand-built
solutions; and (4) the released tool and audit stack, including a headroom
gate that predicts, before any LLM spend, whether a search space will reward
the loop.

\section{Related Work}
LLM-guided program evolution includes AlphaEvolve-style systems and open
engines, with industrial ranking agents (REA, GEARS) closed and DNN-centric.
Automated loss discovery (GLO, AM-LFS, AutoLoss-Zero, CSE-Autoloss) and
hand-derived metric-driven losses (LambdaLoss) motivate our first campaign.
Reflective feedback for LLM-driven optimization (Reflexion, ProTeGi, OPRO,
GEPA) feeds back execution traces and scores; to our knowledge gradient-level
introspection has not previously been used as feedback for loss search. LLM
feature engineering (CAAFE and successors) and the classical target-encoding
leakage literature inform our second campaign. The ICDM 2013 challenge
solutions document the recipes our loop rediscovers. A 2024 survey of
e-commerce LTR notes that datasets with monetary outcome labels remain
scarce, which anchors open multi-objective work to the Expedia benchmark.
None of the above audits discovery against held-out transfer under
loop-realistic noise.

\section{System and Audit Methodology}

\subsection{Background: the evolutionary engine}
%% TODO: add citation for the SkyDiscover framework once its reference is public.
EvoRank builds on SkyDiscover, an open framework for LLM-guided evolutionary
program search in the AlphaEvolve family. The unit of evolution is a source
file in which mutable regions are delimited by EVOLVE-BLOCK markers;
everything outside the markers is frozen, which is how we constrain what the
search may change. Each iteration performs one cycle: sample a parent program
(plus a small set of context programs) from a program database, assemble a
prompt containing their source, scores, and evaluator feedback, make a single
LLM call that proposes an edit to the mutable regions, evaluate the resulting
candidate with a user-supplied \texttt{evaluate(program\_path)} function that
returns a metric dictionary, and insert the scored candidate back into the
database. For readers who know classical evolutionary computation: the LLM
replaces hand-designed mutation and crossover operators, and the prompt
context plays the role of the mating pool.

We use SkyDiscover's AdaEvolve search controller. Its program database is an
island model (two islands here) with periodic ring migration of top programs,
which maintains diversity by evolving semi-isolated subpopulations. Island
selection uses a UCB bandit over decayed improvement rewards, and each island
adapts its exploration intensity from an accumulated improvement signal, so
productive lineages are exploited while stagnating ones explore. In
multi-objective mode the database keeps a Pareto archive: a candidate
survives if no other program dominates it on all objectives (here ndcg,
booking ndcg, and a pooled dollar-weighted revenue metric), with a scalar
fitness key used only for tie-breaking and adaptive-state bookkeeping. On
stagnation the controller can trigger a paradigm step that prompts the LLM
for a high-level strategy change rather than a local edit. All state is
checkpointed every five iterations, so runs resume exactly and every selected
program is preserved with its lineage.

\subsection{Task, data, and evaluators}
Data is the public Expedia Personalized Hotel Search training file (9.9M
impressions, 399k queries), split 70/15/15 by query, with graded labels (5
booked, 1 clicked, 0 otherwise). A fast fitness fold (8k training queries)
keeps candidate evaluation between 5 and 40 seconds; each candidate trains a
gradient-boosted ranker under a fixed budget and is scored on the fold's
validation queries. Degenerate candidates (non-finite gradients or metrics,
timeouts, malformed specifications) receive zero fitness with an explanatory
message that flows back into the next prompt, so guardrails double as
teaching signals.

\subsection{The audit stack}
The audit stack, applied identically to every campaign, is the methodological
core: (a) paired query bootstrap for every comparison, with noise-gated
deltas inside the loop's own feedback; (b) equal-budget random-search
controls; (c) symmetric hyperparameter treatment (post-hoc Optuna tuning of
discovered artifacts at the baseline's budget); (d) a selection-
generalization audit that rescores every checkpoint's selected program on 60k
held-out queries; (e) three-objective hypervolume; and (f) mechanism
ablations of discovered artifacts. Baselines are LambdaMART (the standard
gradient-boosted ranking algorithm, XGBoost rank:ndcg) with 40 trials of
Optuna hyperparameter search, re-tuned separately in every data regime we
report, plus an equal-budget random search over the same candidate space as
the loop.
%% TODO: consider a small figure here summarizing the audit stack.

\section{Campaign 1: Objective Space, a Measured Negative}
The evolvable artifact is a group-aware LambdaMART gradient with access to
raw behavioral signals (relevance grade, booking flag, booking revenue) per
impression; training configuration, features, and metrics are frozen, and the
seed reproduces XGBoost's built-in rank:ndcg within noise.

On the fold the loop selects on, evolved objectives beat the default baseline
and improve steadily. The transfer audit reverses the picture: final selected
objectives gain $+0.0004$ to $+0.0014$ NDCG@10 on held-out queries, the best
evolved objective beats its own seed by $+0.0002$ ($p=0.76$), an equal-budget
random search transfers as well or better ($p=0.03$ in its favor), and every
method is significantly below the tuned baseline ($p=0.0002$). The mechanism
is measured: per-edit effect sizes of 0.002 to 0.005 against a selection-fold
noise floor of $\pm 0.008$ to $0.015$ (1.6k queries). Enlarging the fitness
fold to 6k queries, nearly free since evaluation cost is training-dominated,
halves the noise and triples the mean transferred gain, and is the only
intervention that moved transfer.

Guidance ablations: seeding the prompt with curated domain knowledge reaches
the random-search-budget level on the selection fold in 1 to 3 iterations
versus 7 to 23 without it, and wins three-objective hypervolume in every seed
(mean 252 vs 230, $\times 10^{-6}$), with unchanged transfer. A diagnostic
feedback channel (noise-gated deltas, segment decomposition, a one-step
gradient alignment probe, per-signal gradient-usage probes) is flat across
three seeds on every endpoint. One probe byproduct is worth reporting: the
best discovered revenue term moves only 0.09 percent of gradient mass while
removing it costs 0.009 revenue, so small but systematic gradient biases
compound across boosting rounds.

Mechanism ablation outperformed further search: knocking out components of
the best discovered objective attributed its behavior and revealed that its
per-query adaptive blending was harmful; a constant blend improved all three
metrics at the cost of five evaluations. At full data scale the objective
family is a plateau: all objectives, including default rank:ndcg, land within
0.001 NDCG@10 of each other.

\section{Campaign 2: Pipeline Space, a Measured Positive}
The second campaign evolves two blocks: feature-construction code over
semantically named columns with a leakage-guarded train-fold statistics API,
and a declarative whitelisted pipeline specification (one to three members
from XGBoost, LightGBM, and sklearn families with rank, pointwise, binary, or
campaign-1 custom losses, combined by per-query z-score weighting or rank
averaging). Feedback is stage-attributed so the loop can tell which axis
earned each change: a features line, per-member validation scores, the
ensemble's margin over its best member, noise-gated metric deltas, and
segment decompositions.

A headroom gate precedes any LLM spend: a hand-built candidate from the ICDM
2013 literature must clear 2.5 times the noise floor, and within-query rank
features plus count encodings did ($+0.013$ NDCG@10). The gate surfaced a
three-generation leakage study we report as a practitioner catalog: naive
train-fold target encoding costs $-0.10$ NDCG (the model memorizes its own
labels); leave-one-out encoding is subtly worse at $-0.13$, because
subtracting the row's own label plants a per-row offset perfectly correlated
with that label, which histogram-based trees read as a label detector; and
query-folded out-of-fold encoding is clean. Each generation was caught by the
audit before any search depended on it.

\begin{table}
  \caption{Campaign 2 transfer audit: fast-train protocol, 59{,}902 held-out
  queries. All three seeds beat the tuned baseline; the best seed dominates
  on all three objectives.}
  \label{tab:transfer}
  \begin{tabular}{lcccc}
    \toprule
    Method & Val NDCG@10 & Test NDCG@10 & Test book & Test revenue \\
    \midrule
    Seed pipeline & 0.4218 & 0.4222 & 0.4443 & 0.4105 \\
    LambdaMART + Optuna & 0.4330 & 0.4296 & 0.4527 & 0.4231 \\
    Pipeline s0 & 0.4356 & 0.4316 & 0.4545 & 0.4257 \\
    Pipeline s1 & 0.4359 & 0.4322 & 0.4553 & 0.4215 \\
    Pipeline s2 & 0.4398 & \textbf{0.4352} & \textbf{0.4582} & \textbf{0.4239} \\
    \bottomrule
  \end{tabular}
\end{table}

Three seeds, 50 iterations each, about 150 LLM calls and ten dollars per
seed, independently converged on the same interpretable core: five
within-query percentile ranks (price, star rating, review score, location
scores), count encodings of property and destination identifiers, a lean 20
to 30 column feature set whose program comments cite the measured bloat
evidence, and a diverse three-member ensemble (XGBoost rank:ndcg, LightGBM
lambdarank, plus an extra-trees or binary member) under per-query z-score
weighting. Table~\ref{tab:transfer} shows the transfer audit: roughly 70
percent of selection-fold gains survive on 60k held-out queries, and all
three seeds beat the tuned baseline.

\begin{table}
  \caption{Full-scale comparison (trained on 280k queries, tested on 60k),
  including NDCG@38, the ICDM 2013 challenge metric. The config-equalized
  pipeline is a lower bound (member configuration not re-tuned for the
  pipeline).}
  \label{tab:fullscale}
  \begin{tabular}{lcccc}
    \toprule
    Method & NDCG@10 & NDCG@38 & Book@10 & Revenue@10 \\
    \midrule
    Tuned baseline (fast-regime config) & 0.4349 & 0.4983 & 0.4588 & 0.4300 \\
    Tuned baseline (full-scale config) & 0.4544 & 0.5132 & 0.4805 & 0.4565 \\
    Pipeline s1 (fast-regime config) & 0.4490 & 0.5080 & 0.4738 & 0.4483 \\
    Pipeline s1 (config-equalized) & \textbf{0.4619} & \textbf{0.5185} & \textbf{0.4883} & \textbf{0.4585} \\
    \bottomrule
  \end{tabular}
\end{table}

The full-data regime required its own rigor step and yields a methodology
lesson (Table~\ref{tab:fullscale}): against a baseline tuned in the fast
regime the pipelines win easily, but re-tuning the baseline at full scale
reverses that naive verdict, and only a config-equalized comparison isolates
the discovered structure, which then wins on all four reported numbers.
Relative to the ICDM 2013 challenge report (0.531 NDCG@38 on a hidden,
unreproducible leaderboard, from a 12-model hand-built ensemble trained on
the full data), our numbers are a comparable-protocol approximation on our
own held-out split, in the same ballpark and not claimed as superior.

\section{Discussion}
The two campaigns differ in exactly one variable, the search space, and the
same loop produced noise in one and transfer-audited Pareto dominance in the
other. The deciding quantity is measurable in advance: the ratio of
achievable effect sizes to fitness-fold noise. We suggest a headroom gate as
cheap preregistration for discovery loops, and transfer-audited selection as
the reporting standard. Guidance (prompt knowledge, evidence-rich feedback)
is a speed and interpretability lever, but in our measurements evaluator-side
investment (larger fitness folds, leakage-safe statistics, stage-attributed
rejection messages the model learns from) carried more of the outcome.

Limitations: one dataset, though newer public datasets lack conversion and
revenue labels, which is precisely why open multi-objective LTR remains
anchored to this benchmark; three seeds per condition; NDCG@38 comparisons to
the 2013 challenge are approximate by necessity; revenue is a proxy metric;
and discovered pipelines were audited for transfer, not online effects.

\section{Conclusion}
Publishing an audited failure next to an audited success on the same task is,
we believe, more useful than either alone. The loop is neither magic nor
broken: it is a search procedure whose discoveries are real precisely when
the search space offers effect sizes the fitness signal can resolve. Code,
configs, seeds, per-iteration logs, discovered programs, and audit scripts
are released at \url{https://github.com/shabazpatel/evorank}.

%% TODO: replace with a proper .bib once citations are finalized.
\begin{thebibliography}{9}
\bibitem{liu2013} X. Liu et al. Combination of Diverse Ranking Models for
Personalized Expedia Hotel Searches. arXiv:1311.7679, 2013.
\bibitem{wang2018} X. Wang et al. The LambdaLoss Framework for Ranking Metric
Optimization. CIKM, 2018.
\bibitem{survey2024} A Survey on E-Commerce Learning to Rank.
arXiv:2412.03581, 2024.
\bibitem{reflexion} N. Shinn et al. Reflexion: Language Agents with Verbal
Reinforcement Learning. NeurIPS, 2023.
\bibitem{protegi} R. Pryzant et al. Automatic Prompt Optimization with
Gradient Descent and Beam Search. EMNLP, 2023.
\bibitem{caafe} N. Hollmann et al. Large Language Models for Automated Data
Science: Introducing CAAFE. NeurIPS, 2023.
\bibitem{autolosszero} H. Li et al. AutoLoss-Zero: Searching Loss Functions
from Scratch for Generic Tasks. CVPR, 2022.
\end{thebibliography}

\end{document}
